\documentclass[11pt]{amsart}
\usepackage{amssymb,amsmath,mathtools,booktabs,enumitem}
\usepackage[margin=2.6cm]{geometry}
\usepackage{hyperref}

\newcommand{\stamptime}{\the\numexpr(\time-30)/60\relax:\ifnum\numexpr\time-60*((\time-30)/60)\relax<10 0\fi\the\numexpr\time-60*((\time-30)/60)\relax}
\AddToHook{shipout/foreground}{\put(0.5\paperwidth,-\paperheight+0.55cm){\makebox[0pt]{\tiny\ttfamily unit/doc/\jobname.tex \ --- \ compiled \today, \stamptime}}}

\newtheorem{theorem}{Theorem}[section]
\theoremstyle{definition}
\newtheorem{definition}[theorem]{Definition}
\theoremstyle{remark}
\newtheorem{remark}[theorem]{Remark}

\newcommand{\PSh}{\mathrm{PSh}}
\newcommand{\Sp}{\mathrm{Sp}}
\newcommand{\yo}{\mathbf{y}}
\newcommand{\pres}[1]{\langle #1\rangle}
\newcommand{\Prop}{\mathrm{Prop}}
\newcommand{\tC}{{\textstyle\int}\mathcal C}

\title{The presheaf model of synthetic algebraic geometry:\\ what is going to be checked}
\date{\today}

\begin{document}
\maketitle

\begin{abstract}
We describe the formal verification, now starting, of the main results of
H\"ofer's construction of higher presheaf models (H\"ofer's licentiate
thesis, \emph{Constructive higher sheaf models with applications to
synthetic mathematics}) in the case of synthetic algebraic geometry over a
commutative ring $k$, with the site $\mathcal C$ internal to the base model
as H\"ofer requires. The aim is that nothing is assumed beyond the
parameters of the metatheory. We state what is constructed, what is
proved, in which order, and the points that are still open.
\end{abstract}

\section{Aim}

The model is the presheaf model on a category $\mathcal C$ internal to a
cubical base, relativised to the types that are modal for the cobar
modality $D$ (H\"ofer \S4). In it, the generic $k$-algebra $G$ is to satisfy
the presheaf axioms of synthetic algebraic geometry:

\begin{enumerate}[label=(\arabic*)]
\item \emph{duality} (H\"ofer's Theorem~62): for every finite presentation
$u$, the evaluation map $\pres u \to G^{\Sp(\pres u)}$ is an equivalence;
\item \emph{constancy of $\mathbb N$} (Theorem~61): the constant map
$\mathbb N\to\mathbb N^{\Sp(\pres u)}$ is an equivalence;
\item \emph{projectivity of spectra} (Proposition~76), which uses the
propositional truncation of the model (Proposition~34).
\end{enumerate}

These three statements are what is checked. We do not check that the model
interprets all of type theory. Each axiom is stated directly as a property
of the specific presheaves involved, for instance: the map
$\pres u\to G^{\Sp(\pres u)}$ is an internal equivalence of the model,
uniformly over the context of presentations.

\subsection*{Why $\mathcal C$ must be internal}
Taking $\mathcal C$ to be the ordinary category of finitely presented
$k$-algebras, with constant presheaves on their underlying sets, fails in
general. The constant presheaf on a set quotient $L[X]/q$ is not the
$h$-set quotient of the model, as shown in the erratum of Coquand, H\"ofer
and Sattler to the paper on synthetic algebraic geometry. With
$\mathcal C$ internal to the base, set quotients of the metatheory are
never used: all quotients are formed in the base model.

\section{What is assumed}

Nothing beyond the parameters of the metatheory. The metatheory is Unit
(\texttt{github.com/coquand/unit}), with function extensionality,
propositional univalence and the univalence and replacement parameters of
its developments. The cubical base is the poset model, built inside Unit.
Every construction in the model is a definition, and every property a
proof.

Two things fall outside the checker by nature:
\begin{enumerate}[label=(\alph*)]
\item the agreement between each concrete statement in the model and the
axiom of synthetic algebraic geometry it interprets. There is no general
soundness theorem, so this is checked by reading the statement;
\item the equivalence between our site and H\"ofer's
(Remark~\ref{rem:site}). This is an argument in the model, and it is not
used by any proof.
\end{enumerate}

\section{The construction}

\subsection{Tools in the base}
\begin{enumerate}[label=(B\arabic*)]
\item \emph{Propositional truncation}, as in the paper on synthetic
algebraic geometry, \S8.3. An element of $\|T\|$ is either $\mathrm{inc}(a)$,
or $\mathrm{ext}(u,\psi,v)$ with $\mathrm{ext}(u,\psi,v)=v$ on the extent
$\psi$. It is an ordinary inductive family of normal forms, with no
interval and no higher inductive type. Most of it is already in the core
library; what remains is composition and stability under substitution.
\item \emph{The universe of propositions} $\Prop$ of the base, from its
univalent universe.
\item \emph{$h$-set quotients} $A/R$, by Voevodsky's construction: the
image of $a\mapsto R(a,-)$ in $A\to\Prop$. Without replacement in the base,
$A/R$ lies one universe level above $A$. Quotients are never put inside a
universe, so this level shift suffices.
\item \emph{The algebra of the base}, for models in $h$-sets of the base:
\begin{itemize}
\item free $k$-algebras $F_n$, as quotients of coefficient lists;
\item presented algebras, each one quotient of terms by the congruence
generated by the theory and the relations;
\item relative presentations, and H\"ofer's Lemmas~42 and 50 (pushouts).
\end{itemize}
These results are already proved in the metatheory (Tier~1). They must be
proved again for $h$-sets of the base, since the base is a model built
inside the metatheory.
\end{enumerate}

\subsection{The site}
\begin{definition}
The internal category $\mathcal C^{\mathrm{op}}$ has as objects the finite
presentations $(n,m,s,t)$ with $s,t\colon \mathrm{Fin}\,m\to F_n$. Its
morphisms are H\"ofer's (Construction~39) between the presented algebras,
and composition is strict. Its objects and morphisms are types of the base,
so $\mathcal C$ varies over the cubes: $F_n$ is a quotient formed in the
base.
\end{definition}

\begin{remark}\label{rem:site}
This category is equivalent to H\"ofer's classifying category of finitely
presentable algebras, by the argument of his Lemma~40. Taking
presentations as objects avoids a universe of models, and the smallness
argument of his Remark~38, which is the only place where H\"ofer uses
replacement.
\end{remark}

The model is taken over the Grothendieck construction $\tC$, a site
with interval and cofibrations pulled back from the base. The fibrancy
theory of the base is generic in its site, so at $\tC$ it gives:
\begin{itemize}
\item fibrant types with $\Sigma$, $\Pi$ and identity types;
\item the universe and the truncation;
\item contractibility and homotopy facts.
\end{itemize}
Levelwise notions are fibrancy at the slices of the base.

\subsection{The cobar modality and H\"ofer's \S4}
Over $\tC$ we establish:
\begin{itemize}
\item $D$, the cobar modality (Definition~10), over chains of $\mathcal C$
at each stage;
\item Lemmas~12, 14 and 16, Proposition~13 and Corollary~15;
\item Proposition~17: for $D$-modal types, a levelwise equivalence is an
equivalence;
\item Theorem~29 for the identity modality: $D$ is a descent data
operation.
\end{itemize}
These results are already checked when $\mathcal C$ is constant. The
proofs stay the same; the indices now move with the cubes.

\subsection{The objects and the axioms}
\begin{itemize}
\item $G := D(\yo F_1)$ (Construction~45). Its ring structure comes from
the elements of $F_n$ (Construction~44), and its fibre at $x$ is the
algebra of $x$.
\item $\pres u$: the $D$-modal reflection of the quotient, at $\tC$, of
the terms over $G$. Its fibres are the presented algebras, by Lemma~52.
\item $\Sp(\pres u)$ and $G^{\Sp(\pres u)}$, formed with $\Sigma$, $\Pi$
and identity types. Their fibres come from Lemmas~54 and 59 and
Proposition~58 (Constructions~55--57).
\item $\mathbb N_D := D(\Delta\mathbb N)$.
\end{itemize}
The axioms then follow as in H\"ofer's \S6. For duality: $\pres u$ and
$G^{\Sp(\pres u)}$ are $D$-modal, so by Proposition~17 it is enough that
evaluation is a levelwise equivalence. This is the algebra of the base
read at the fibres, which is already checked in the metatheory and is
proved again in the base. Projectivity uses Proposition~34, with the
truncation of the model defined from (B1) through the right adjoint $R$,
and H\"ofer's Proposition~76.

\section{Order, sizes, open points}

\begin{center}
\begin{tabular}{@{}llr@{}}
\toprule
Step & Content & Lines (estimate)\\
\midrule
0 & base tools (B1)--(B3), constant types, lists & 1{,}300--2{,}000\\
1 & algebra of the base (B4) & 6{,}000--11{,}000\\
2 & the internal category $\mathcal C$ & 600--800\\
3 & the sites $\tC$ and the slices, pullbacks & 1{,}000--1{,}700\\
4 & $D$, Lemmas 12--16, Props.~13, 17, Thm.~29, Lemma 54 & 5{,}500--7{,}000\\
5 & $G$, $\pres u$, $\Sp$, $G^{\Sp}$, Prop.~58, duality & 3{,}100--4{,}600\\
6 & constancy of $\mathbb N$, projectivity (Prop.~34) & 2{,}000--3{,}000\\
\midrule
& total & 19{,}500--30{,}000\\
\bottomrule
\end{tabular}
\end{center}

Step~1 decides the size. A pilot comes first: the free algebra $F_1$ as a
quotient, with its universal property into an $h$-set ring of the base. It
measures how much longer the base proofs are than their counterparts in
the metatheory.

\medskip
\noindent\textbf{Open points.}
\begin{enumerate}[label=(\roman*)]
\item \emph{Universe levels.} The base has one universe step. $F_n$, and
with it $\mathcal C$, $G$ and the whole model, live one level up. The
quotient $\pres u$ at $\tC$ is a quotient of terms over $G$, whose
congruence takes values in identity types of $G$. Whether these are
propositions of the available universe, or whether a second universe step
is needed, is being settled.
\item \emph{Proposition~34 and the truncation of the model}: built from
(B1) by the right adjoint $R$, with no assumption; it remains to be sized
over $\tC$.
\end{enumerate}

\end{document}
